\section{Statement of the Problem}

This study examines the main limitation of the current SAGE approach: its reliance on the pixel2style2pixel (pSp) encoder for GAN inversion. The pSp-based inversion stage can yield entangled latent codes and reduce control over semantic edits. In few-shot image generation, this limitation affects the production of semantically consistent and high-fidelity images from limited reference examples. Loss of local detail during inversion can also reduce the usefulness of generated images for downstream tasks because reconstructed images may no longer preserve the information present in the original samples.

The study therefore investigates whether improving the inversion and feature-editing stages of SAGE can strengthen both generated-image quality and editability. Specifically, the study seeks to answer the following research questions:

\begin{enumerate}
    \item Can integrating the StyleFeatureEditor inverter and Feature Editor improve latent inversion fidelity and preserve image details for style-based editing within the SAGE framework?

    \item Do these modifications influence the consistency, diversity, and quality of images generated in the few-shot setting, particularly for novel classes in the 102 Flowers and Animal Faces datasets?

    \item How well does FeatureSAGE perform in terms of editing reliability and disentanglement in few-shot scenarios?

    \item How does FeatureSAGE compare with existing editing-based few-shot image generation methods, specifically AGE and SAGE, in terms of edit stability, category or identity consistency, and overall image quality?
\end{enumerate}
